\Needspace{12\baselineskip}
\originalchapter{Lab3：金属、堆垛与铁电的全部公开任务}
原Lab3截止于2005年4月5日，使用PZ-LDA超软赝势。它把上一章收敛方法推进到小相能差、金属自旋、堆垛缺陷与内部位移。Co的指定命名赝势可从当前QE官方库取得，本项目实际运行对应输入族；公开Lab3归档内也确有原始 \texttt{Ba.UPF,Ti.UPF,O.UPF}，本次已安全解压并核验哈希；原程序版本与当前运行版本仍需区分。本章保留全部问题及官方结果，同时明确区分原结果、现代Co复算与原归档输入的实际任务结果。

\originalsection{Co的晶格和相能差}
\begin{exercise}\label{task:lab3-1}
（原题1，五项）用PZ-LDA，波函数/密度截断30/250 Ry起步，自己选择k网格，能差目标约1 meV/atom。
(1)画FCC Co能量随晶格常数变化，用公开元素资料取起始值。
(2)HCP先固定理想 $c/a=\sqrt{8/3}$ 找a，再固定a改变c；在附近若干a分别画c扫描，找到二维最低点，比较实验a与c/a，不必采用小于.01 bohr的过细网格。
(3)比较HCP/FCC每原子能量，判断稳定相并猜测加压、升温是否可能转相；HCP胞含2原子、FCC原胞含1而非4。
(4)A、B层平面位置已给 $(1/3,2/3),(2/3,1/3)$，求C；把ABC的FCC构造成三层六角胞，求立方与六角a的关系，用理想三层c/a核对相能差。
(5)在HCP中引入 $AB|CACACAC\cdots$ 堆垛缺陷，固定a、c/a与原子位置，考虑周期边界，画额外能量随层数变化；不超过10层，平面用 $8\times8$ k网格，c方向网格随胞长反比调整。
\end{exercise}
\begin{solution}
金属费米面使k网格与展宽共同影响小能差。原手册参数为
\begin{center}\begin{tabular}{ll}
\texttt{nspin}&2\\\texttt{starting\_magnetization}&.7\\
\texttt{smearing}&\texttt{'cold'}\\\texttt{degauss}&.03 Ry\\
波函数/密度截断&30/250 Ry
\end{tabular}\end{center}
这些输入比只写“DFT”更能决定可复核结果。原官方收敛图建议FCC $16^3$，HCP $12\times12\times6$，还单独检验c方向的网格。新程序保留原偏移 $(0,0,0)$，分别作这两类收敛扫描，不能将绝缘钻石的 $8^3$ 门槛直接移给Co。

原官方结果：FCC $a\simeq6.483$ bohr；HCP二维扫描在 $a\simeq4.58,c\simeq7.438$ bohr附近最低，$c/a\simeq1.624$。原答案给出的LDA/实验比较（\AA）是FCC 3.43/3.54，HCP a为2.42/2.51、c为3.94/4.07、c/a为1.624/1.623。必须把低温理论状态与实验温压条件分开；比例较准确可能含误差相消，不证明两个长度各自无误差。

原官方每原子能量为FCC $-74.22270$ Ry、HCP $-74.22350$ Ry，HCP低约11 meV/atom；每原子体积为68.12与67.56 bohr$^3$。在这些近零压状态附近，$\Delta H\simeq\Delta E+P\Delta V$ 的小压展开使较小体积HCP继续有利；升温的稳定性还需振动、磁性和电子熵，不能由“体积较大”直接证明熵较大。原答案只作温度转相的粗略猜测，其限制应保留。

C层为 $(0,0)$。六角平面a是FCC最近邻距，因此 $a_{\rm hex}=a_{\rm cub}/\sqrt2$，三层高度给 $c/a=3\sqrt{8/3}/2=\sqrt6$。取原FCC a给 $a_{\rm hex}\simeq4.584$ bohr。三层原子z可取0、1/3、2/3，整体沿z平移不改能量。原官方用 $12\times12\times4$ 网格得 $-74.22259$ Ry/atom，与立方FCC结果约差1.5 meV/atom；这量级需继续检查k积分，不能把几种胞总能量不经原子数归一化直接相减。

堆垛缺陷的周期序列应具体列出，例如3、5、7、9层为ABC、ABCAC、ABCACAC、ABCACACAC；每添两层恢复较远处HCP环境。令 $E_n$ 为n层胞能、$E_2$ 为两层HCP参考，缺陷超额为 $\Delta E_n=E_n-nE_2/2$，若归每原子则再除n。两者不能混用：总超额趋常数时，每原子超额本来就趋零。是否有一个或多个周期界面须由所建序列检查；换算面积能需按实际缺陷数和面积归一化。原官方图在9层仍未显示充分收敛，原题又限制不超过10层，故不能把最大n点宣布为无限层值。

当前Co程序实际保存FCC、HCP二维、六角FCC和这些缺陷序列的输入与SCF数据；具体现代结果和收敛状态另有独立CSV/QA，与上述官方数值分别标记。它没有把其他赝势、不同磁态或不同胞数的计算混成同一个参考。
\end{solution}

\originalsection{BaTiO$_3$的对称性破缺}
\begin{exercise}\label{task:lab3-2}
（原题2，四项）(1)画立方BaTiO$_3$能量随a变化，用 $4^3$、$(1,1,1)$ 偏移网格，精细扫描求a；室温实验约4 \AA{}。(2)固定该a，沿一立方轴位移Ti，每个Ti位移下使O充分弛豫，找能量最低的Ti位移并与立方态比较。(3)以该构型开始，同时弛豫Ti和O，报告最终原子坐标与能量。(4)原题末项写“PbTiO$_3$哪一相更稳定及如何联系铁电性”，而标题、输入和官方答案均为BaTiO$_3$；保留此化学式不一致，不擅自生成另一种材料的数据。
\end{exercise}
\begin{solution}
原题括注“1 ang=.529177 bohr”的换算方向误写；正确是1 bohr约.529177 \AA{}，所以4 \AA{}约7.56 bohr。扫描时a、坐标单位和原胞原子数须统一。Ba位于立方角、Ti近体心、O在三个面心，立方约束使宏观极化按对称性消失；允许内部位移后再比较低能构型，才能讨论位移型铁电机制。

原官方(a)给 $a\simeq7.52$ bohr，即3.979 \AA{}。(b)文字给最低Ti位移约 $0.08a_0$，能差 $3.726\times10^{-4}$ Ry；该处 $a_0$ 的单位须与实际输入/图轴核对，不能将字面“.08”直接拼成alat分数坐标。开放Ti与O后，官方最终alat坐标为
\[
\begin{array}{c|rrr}
\mathrm{Ba}&0&0&0\\
\mathrm{Ti}&.5&.5&.507394883\\
\mathrm{O}_1&.5&.5&-.019383444\\
\mathrm{O}_2&.5&0&.485432812\\
\mathrm{O}_3&0&.5&.485432812
\end{array}
\]
官方总能量为 $-303.6868329953$ Ry，解释该能量与前项在其SCF准则内基本相同。这些是\emph{原公开答案}，不是本项目重新执行BaTiO$_3$的结果。负O坐标按周期边界可加1等价，但原坐标先原样保存，以便核对相对位移。

复用方法是：先收敛立方a；每个固定Ti位移只开放O的力优化，保存力阈值、最大力和能量；再释放Ti，避免把约束优化点和全弛豫点混为一谈。以极化方向反转的对称构型检查双势阱和能量偶性，但小单胞不包含畴与有限温相变。带电离子的相对位移可产生偶极矩，实际晶体极化的定量计算还涉及Born有效电荷或Berry相，而非只用名义离子电荷。

官方末项回答BaTiO$_3$的破缺对称低能态，不能作为PbTiO$_3$的原计算证据。若按原题Pb字样执行，需要Pb的独立赝势及对应输入，再做相同的约束/弛豫比较。本次使用公开归档中的原Ba/Ti/O文件，而不是按通用名随意选现代赝势；归档与各输入哈希保留，实际复算结果与以上官方答案分别标记。
\end{solution}

\originalsection{Co高压磁性与相稳定}
\begin{exercise}\label{task:lab3-extra}
（原可选题）比较HCP Co磁矩与实验；对HCP/FCC画磁矩随压强变化、总能量随压强变化，判断高压转相及与题1的粗略猜测是否一致；可参考原题列出的Yoo等2000年论文，但不要求复制第三方原图。
\end{exercise}
\begin{solution}
原公开答案给HCP约1.59 $\mu_B$/atom、实验1.72，并报告两种结构在高压转为非磁、约100 GPa附近非磁FCC更有利。这是官方计算结论的保留，不是用题1近零压的两点体积就能证明的结果。压强会改变磁态和能量曲线，因而可使低压外推失效。

正确稳定性比较在给定P下用 $H(P)=\min_{V,\mathrm{shape},\mathrm{magnetic\ branch}}[E+PV]$，而不是把不同压力下的E直接相减。体积扫描的 $P=-\dd E/\dd V$ 或SCF平均应力提供P，HCP还应检查c/a与应力各向异性。磁性初值可能收敛到不同局部解，需要多个初磁矩或非磁分支。

当前程序实际扫描两结构的压缩状态、保留.7与0两种初磁矩，输出平均应力和磁矩。此扫描固定规定形状/c/a，属于原任务的可执行磁分支与压力诊断；不能冒充已完成全部可变胞焓优化或严格定位100 GPa相界。后者的可复用方法与官方结果仍完整列出，原第三方论文图未复制。
\end{solution}
